% -*- coding: utf-8 -*-
% !TEX program = xelatex
\documentclass[plain,noanswer,medmath]{../../template/jnuexam}
\usepackage{../../template/kysx}

\begin{document}


\examtitle{1995}{5}


\loadproblems{4}{1995P4}%载入试卷四


\makepart{填空题}{本题共5小题，每小题3分，满分15分}


\begin{problem}
设$\lim_{x \to \infty} \left(\frac{1 + x}{x} \right)^{ax} = \int_{-\infty}^a t\e^t \dt$，
则常数$a$等于 \fillin{}．
\end{problem}


\useproblem{4}{1}{2}%【同试卷四第一(2)题】


\useproblem{4}{1}{3}%【同试卷四第一(3)题】


\useproblem{4}{1}{4}%【同试卷四第一(4)题】


\begin{problem}
设$X$是一个随机变量，其概率密度为$f(x) =\begin{cases}
 1 + x, & -1 \le x \le 0, \\
 1 - x, & 0 < x \le 1, \\
     0, & \text{其他},
\end{cases}$则方差$DX =$ \fillin{}．
\end{problem}


\makepart{选择题}{本题共5小题，每小题3分，满分15分}


\useproblem{4}{2}{1}%【同试卷四第二(1)题】


\useproblem{4}{2}{2}%【同试卷四第二(2)题】


\begin{problem}
设$n$维行向量$\alpha = \left(\frac{1}{2},0,0,\frac{1}{2} \right)$，
矩阵$A = E - \alpha^T\alpha$, $B = E + 2\alpha^T\alpha$，
其中$E$为$n$阶单位矩阵，则$AB$等于\pickout{}
\begin{abcd}
\item $O$．		
\item $- E$．		
\item $E$．		
\item $E + \alpha^T\alpha$．
\end{abcd}
\end{problem}


\begin{problem}
设矩阵$A_{m \times n}$的秩为$r(A) = m < n$，$E_m$为$m$阶单位矩阵，下述结论中正确的是\pickout{}
\begin{abcd}
\item $A$的任意$m$个列向量必线性无关．
\item $A$的任意一个$m$阶子式不等于零．
\item 非齐次线性方程组$Ax = b$一定有无穷多组解．
\item $A$通过初等行变换，必可以化为$(E_m,0)$的形式．
\end{abcd}
\end{problem}


\useproblem{4}{2}{5}%【同试卷四第二(5)题】


\makepart{}{本题满分6分}%【同试卷四第三题】

\useproblem{4}{3}{0}


\makepart{}{本题满分6分}

\begin{problem}
求不定积分$\int (\arcsin x)^2 \dx$．
\end{problem}


\makepart{}{本题满分7分}%【同试卷四第八题】

\useproblem{4}{8}{0}


\makepart{}{本题满分6分}%【同试卷四第七题】

\useproblem{4}{7}{0}


\makepart{}{本题满分5分}

\begin{problem}
设$f(x)$在区间$[a,b]$上连续，在$(a,b)$内可导，证明：在$(a,b)$内至少存在一点$\xi$，使得
$$\frac{bf(b) - af(a)}{b - a} = f(\xi) + \xi f'(\xi).$$
\end{problem}


\makepart{}{本题满分9分}

\begin{problem}
求二元函数$z = f(x,y) = x^2 y(4 - x - y)$在由直线$x + y = 6$，$x$轴和$y$
轴所围成的闭区域$D$上的极值、最大值与最小值．
\end{problem}


\makepart{}{本题满分8分}

\begin{problem}
对于线性方程组$\begin{cases}
  \lambda x_1 + x_2 + x_3 = \lambda - 3, \\
  x_1 + \lambda x_2 + x_3 = - 2, \\
  x_1 + x_2 + \lambda x_3 = - 2,
\end{cases}$讨论$\lambda$取何值时，
方程组无解、有唯一解和有无穷解？在方程组有无穷解时，试用导出组的基础解系表示全部解．
\end{problem}


\makepart{}{本题满分8分}

\begin{problem}
设三阶矩阵$A$满足$A\alpha_i = i\alpha_i$($i = 1,2,3$)，
其中列向量$\alpha_1 = (1,2,2)^T$, $\alpha_2 = (2,-2,1)^T$, $\alpha_3 = (-2,-1,2)^T$，
试求矩阵$A$．
\end{problem}


\makepart{}{本题满分8分}%【同试卷四第十一题】

\useproblem{4}{11}{0}


\makepart{}{本题满分7分}

\begin{problem}
假设随机变量$X$服从参数为$2$的指数分布，证明：$Y = 1 - \e^{-2X}$在区间$(0,1)$上服从均匀分布．
\end{problem}



\end{document}
